% ==============================================================================
% ACTIVIDAD S6 - REDACCION EN CONTEXTOS VIRTUALES / UnADM
% Producto: Lista de verificacion para evaluar un texto digital
% Archivo BibTeX asociado: Actividad_S6_Lista_Verificacion_Texto_Digital.bib
% ==============================================================================

\documentclass[12pt,letterpaper,oneside]{article}

% -------------------- PAQUETES --------------------
\usepackage[spanish, es-tabla]{babel}
\usepackage[utf8]{inputenc}
\usepackage[T1]{fontenc}
\usepackage{geometry}
\usepackage{setspace}
\usepackage{helvet}
\usepackage{array}
\usepackage{longtable}
\usepackage{booktabs}
\usepackage{enumitem}
\usepackage{fancyhdr}
\usepackage{hyperref}
\usepackage{natbib}
\usepackage{xcolor}
\usepackage{amssymb}
\usepackage{titlesec}
\usepackage{ragged2e}

\renewcommand{\familydefault}{\sfdefault}
\geometry{top=2.2cm,bottom=2.2cm,left=2.3cm,right=2.3cm}
\onehalfspacing
\setlength{\parindent}{0pt}
\setlength{\parskip}{6pt}
\setcitestyle{authoryear,open={(},close={)}}

\hypersetup{
  colorlinks=true,
  linkcolor=black,
  citecolor=black,
  urlcolor=black,
  pdftitle={Lista de verificacion para evaluar textos digitales},
  pdfauthor={Martin Jonathan de la Cruz}
}

\pagestyle{fancy}
\fancyhf{}
\lhead{Redaccion en contextos virtuales}
\rhead{Semana 6}
\cfoot{\thepage}

\titleformat{\section}{\large\bfseries}{\thesection.}{0.5em}{}
\titleformat{\subsection}{\normalsize\bfseries}{\thesubsection.}{0.5em}{}

\newcommand{\checkyes}{\textbf{\checkmark}}
\newcommand{\checkpartial}{\textbf{Parcial}}
\newcommand{\checkboxno}{\textbf{No}}

\begin{document}

% -------------------- PORTADA --------------------
\begin{titlepage}
\centering
\vspace*{1cm}
{\Large \textbf{Universidad Abierta y a Distancia de Mexico}}\\[0.4cm]
{\large Licenciatura en Derecho}\\[1.2cm]

{\LARGE \textbf{Lista de verificacion para evaluar textos digitales}}\\[0.3cm]
{\large \textbf{Actividad Semana 6}}\\[1.2cm]

\begin{tabular}{rl}
\textbf{Asignatura:} & Redaccion en contextos virtuales \\
\textbf{Tema:} & Aspectos relevantes para escribir en formatos digitales \\
\textbf{Alumno:} & Martin Jonathan de la Cruz \\
\textbf{Matricula:} & ES2611202040 \\
\textbf{Docente:} & Ana Rosa Arceo Luna \\
\textbf{Fecha:} & \today \\
\end{tabular}

\vfill
\textbf{Texto evaluado:} \emph{El Estado mexicano en la Constitucion Politica de los Estados Unidos Mexicanos}.\\
\textbf{Producto:} lista de verificacion completa y analisis breve.
\end{titlepage}

% -------------------- RESUMEN --------------------
\section*{Resumen}
Esta actividad evalua un texto digital de contenido juridico mediante una lista de verificacion centrada en proposito, claridad, organizacion, tono, legibilidad y correccion. El criterio de lectura no se limita a decir si el texto ``esta bien'', sino a ubicar que tan funcional resulta para comunicar una idea academica en pantalla. La postura que guia esta revision es directa: un texto digital puede tener informacion correcta, pero si no jerarquiza, no cita con precision y no orienta al lector, pierde fuerza academica y se vuelve solo exposicion general \citep{nunez2015,correa2018}.

\section{Texto digital seleccionado}
El texto seleccionado es un estudio de caso titulado \emph{El Estado mexicano en la Constitucion Politica de los Estados Unidos Mexicanos}. Su proposito es explicar la configuracion juridica del Estado mexicano a partir de los articulos 1, 39, 40, 41, 42, 49 y 133 constitucionales. El texto sostiene que Mexico se organiza como una republica representativa, democratica, laica y federal; ademas, vincula la soberania popular, la division de poderes, el territorio, los derechos humanos y la supremacia constitucional como elementos estructurales del Estado \citep{constitucion2026}.

Desde el contenido juridico, el texto es util porque parte de normas constitucionales concretas. Esta base conecta con una idea central de la teoria juridica: la norma es la que otorga significado juridico a ciertos actos, instituciones y relaciones \citep{kelsen1982}. Tambien permite leer al Estado como estructura de poder, orden y justicia, no solo como definicion politica \citep{paredes2020,rojas2018}. Sin embargo, desde la escritura digital, el punto critico no es solo tener informacion correcta, sino convertirla en una lectura ordenada, escaneable y con cierre analitico. En textos academicos digitales, la claridad depende de la organizacion del argumento, la coherencia entre parrafos y el uso visible de fuentes \citep{corte2019}.

\section{Lista de verificacion}

\begingroup
\small
\setlength{\tabcolsep}{5pt}
\renewcommand{\arraystretch}{1.35}
\begin{longtable}{@{}p{0.19\linewidth}p{0.49\linewidth}p{0.10\linewidth}p{0.18\linewidth}@{}}
\caption{Evaluacion del texto digital seleccionado}\label{tab:lista-verificacion}\\
\toprule
\textbf{Categoria} & \textbf{Criterio} & \textbf{Cumple} & \textbf{Observacion breve} \\
\midrule
\endfirsthead
\toprule
\textbf{Categoria} & \textbf{Criterio} & \textbf{Cumple} & \textbf{Observacion breve} \\
\midrule
\endhead
\bottomrule
\endlastfoot

\textbf{Proposito y claridad} & El texto tiene un proposito claro. & \checkyes & Explica el fundamento constitucional del Estado mexicano. \\
\midrule
\textbf{Proposito y claridad} & La idea principal se identifica facilmente desde el inicio. & \checkyes & La tesis aparece desde el primer parrafo. \\
\midrule
\textbf{Organizacion y estructura} & Los parrafos son breves y estan bien delimitados. & \checkyes & La division por articulos facilita la lectura. \\
\midrule
\textbf{Organizacion y estructura} & La informacion se presenta de forma logica y ordenada. & \checkyes & Avanza de soberania, forma de Estado y poderes hacia territorio, derechos y supremacia. \\
\midrule
\textbf{Organizacion y estructura} & Se utilizan conectores que dan coherencia al texto. & \checkpartial & Hay conectores basicos, pero podria mejorar la transicion entre bloques. \\
\midrule
\textbf{Lenguaje y tono} & El lenguaje es adecuado al contexto academico o profesional. & \checkyes & Utiliza vocabulario juridico formal. \\
\midrule
\textbf{Lenguaje y tono} & Se evita el uso de expresiones coloquiales o informales. & \checkyes & No presenta expresiones informales. \\
\midrule
\textbf{Lenguaje y tono} & El tono es claro, formal y respetuoso. & \checkyes & Mantiene neutralidad academica. \\
\midrule
\textbf{Legibilidad digital} & El texto es facil de leer en pantalla. & \checkpartial & Es claro, pero requiere subtitulos o marcadores visuales para escanear mejor. \\
\midrule
\textbf{Legibilidad digital} & No hay parrafos excesivamente largos. & \checkyes & Los parrafos son manejables. \\
\midrule
\textbf{Legibilidad digital} & La estructura facilita la lectura rapida o escaneabilidad. & \checkpartial & La lectura mejoraria con viñetas, subtitulos o tabla normativa. \\
\midrule
\textbf{Correccion} & No presenta errores ortograficos. & \checkyes & No se detectan errores relevantes. \\
\midrule
\textbf{Correccion} & No presenta errores gramaticales. & \checkyes & La sintaxis es comprensible y formal. \\
\midrule
\textbf{Fuentes y respaldo} & Integra correctamente las fuentes consultadas. & \checkpartial & Cita la Constitucion al final, pero faltan citas dentro del cuerpo. \\
\midrule
\textbf{Analisis academico} & Presenta postura, interpretacion o valoracion propia. & \checkpartial & Cierra con sintesis, pero el analisis personal podria ser mas visible. \\

\end{longtable}
\endgroup

\section{Analisis breve del texto evaluado}
El texto cumple su funcion principal: explicar como la Constitucion organiza al Estado mexicano a partir de soberania, federalismo, division de poderes, territorio, derechos humanos y supremacia constitucional. Su mayor fortaleza es la secuencia: cada articulo cumple una funcion dentro del sistema y eso permite entender la estructura estatal sin perder el hilo. Tambien mantiene un tono academico y formal, lo cual lo hace adecuado para un entorno universitario digital. La debilidad esta en que se queda muy expositivo: menciona los articulos, pero no siempre explica su peso juridico o su consecuencia practica. Ademas, la fuente aparece concentrada al final, cuando lo mas correcto seria integrar citas dentro del cuerpo para sostener cada bloque normativo y distinguir entre derecho positivo, legitimidad y derechos humanos \citep{finnis2017}. Como ajustes concretos, agregaria subtitulos breves por eje: soberania, forma de Estado, poder publico, territorio y supremacia constitucional. Tambien incorporaria una tabla final que relacione articulo, funcion y efecto juridico. Asi el texto no solo informa; orienta, jerarquiza y permite leer el sistema constitucional con mayor criterio.

\section{Ajustes propuestos}
\begin{enumerate}[label=\arabic*.]
    \item Integrar citas parenteticas dentro del texto, por ejemplo, despues de explicar cada articulo constitucional, para evitar que la referencia quede como adorno final.
    \item Añadir subtitulos o una tabla de sintesis que mejore la lectura en pantalla y permita ubicar rapidamente la funcion de cada articulo.
    \item Incluir una postura de cierre mas fuerte: no solo decir que el Estado mexicano esta estructurado constitucionalmente, sino explicar que esa estructura opera como limite al poder y como garantia minima para la ciudadania.
\end{enumerate}

\section{Conclusion}
La evaluacion muestra que el texto seleccionado tiene base juridica, orden y tono academico, pero necesita mayor trabajo de escritura digital. En frio: no basta con poner informacion correcta; hay que diseñar la lectura. Un texto academico en formato digital debe funcionar como sistema de orientacion: abrir con una idea clara, organizar por niveles, citar donde corresponde y cerrar con criterio propio. Esa mejora no cambia el contenido constitucional, pero si cambia su fuerza comunicativa.

\bibliographystyle{apalike}
\bibliography{Actividad_S6_Lista_Verificacion_Texto_Digital}

\end{document}
